%% kapa: exam=1 round=21 year=2010 period=1 form=A subject=civil_law subject_ko=민법 seq=40 number=40
%% source: Q-Net 감정평가사 기출문제 (cst00309) artlSeq=5024181 file=extracted/1008/5024181/2010년 제21회 감정평가사 1차시험 기출문제/1교시.hwp
%% conversion: kordoc 4.24.51 -> tools/gen1.py (auto); review=proofread-sonnet; flags=none
\begin{question}{40}
甲이 乙에게 대여한 1억원의 채권을 담보하기 위해 乙소유의 X토지 위에 저당권이 설정되었다. 다음 설명으로 옳지 않은 것은? (다툼이 있으면 판례에 의함)
\begin{choices}
  \item 甲의 저당권은 원본, 이자, 위약금, 채무불이행으로 인한 손해배상 뿐만 아니라 저당권의 실행비용도 담보하는 것이 원칙이다.
  \item 乙의 고의나 과실로 인하여 X토지의 가액이 현저히 감소된 경우, 甲은 乙에게 그 원상회복 또는 상당한 담보제공을 청구할 수 있다.
  \item 저당권이 설정된 후 乙이 X토지에 건물을 축조하여 소유하고 있는 경우, 甲은 X토지와 함께 그 건물에 대해서도 경매를 청구할 수 있으며, 이 경우 그 건물의 매각대금으로부터 우선변제를 받을 수 있다.
  \item 甲이 저당권부 채권을 丙에게 양도하려면 채권양도의 요건 이외에 저당권 이전에 대한 물권적 합의와 등기가 있어야 하는데, 여기서 물권적 합의는 甲ㆍ丙간에만 있으면 충분하고 그 외에 乙에게까지 있어야 하는 것은 아니다.
  \item 1억원의 대여금채권이 소멸한 후에 甲이 별도로 취득한 乙에 대한 채권을 담보하기 위해서 종전의 저당권등기를 유용하기로 합의한 경우, 그 합의 이전에 丁명의의 지상권이 설정되어 있다면 丁에게 저당권의 유효를 주장할 수 없다.
\end{choices}
\end{question}
\answer{3}
